\documentclass[10pt,a4paper]{amsart}
\usepackage[margin=19mm]{geometry}
\usepackage{amsmath,amssymb,mathtools}
\usepackage[scr=rsfs,cal=euler]{mathalfa}
\usepackage{xfrac}
\usepackage[unicode,hidelinks,bookmarks=false]{hyperref}
\newcommand{\ie}{i.e.\ }
\newcommand{\cf}{cf.\ }
\newcommand{\resp}{respectively\ }
\newcommand{\Subsection}{\subsection}
\providecommand{\nobreakdash}{\nobreak\hbox{-}}
\setlength{\emergencystretch}{2em}
\multlinegap0pt
\usepackage{bm}
\usepackage{bbm}
\usepackage{dsfont}
\usepackage{xspace}
\usepackage{cancel}
\usepackage{mathtools}
\usepackage[shortlabels]{enumitem}
\usepackage{amsthm}
\makeatletter
\@ifundefined{theorem}{\newtheorem{theorem}{Theorem}}{}
\@ifundefined{corollary}{\newtheorem{corollary}[theorem]{Corollary}}{}
\@ifundefined{lemma}{\newtheorem{lemma}[theorem]{Lemma}}{}
\@ifundefined{proposition}{\newtheorem{proposition}[theorem]{Proposition}}{}
\makeatother
\title[Forbidden Four Cycle, Star Graphs and Isometric Embeddings]{Forbidden Four Cycle, Star Graphs and Isometric Embeddings}
\author{Oleksiy Dovgoshey, Olga Rovenska}
\date{Source: arXiv:2510.01667v1 (CC BY 4.0)}
\begin{document}
\maketitle

\noindent\emph{Source and licence.} Forbidden Four Cycle, Star Graphs and Isometric Embeddings. Version arXiv:2510.01667v1 (CC BY 4.0), distributed under the Creative Commons Attribution 4.0 International licence, as recorded in the arXiv licence metadata for this exact version and shown on the abstract page https://arxiv.org/abs/2510.01667v1. Full rights notice: licenses/arxiv-2510.01667-CC-BY-4.0.txt.\par\smallskip

\noindent\textit{Editorial scope.} Counted unit: the Lemma whose source label is \texttt{476} in the article (source0.tex lines 1454--1464, source label \texttt{476}), delivered together with the article's own complete original proof of it (source0.tex lines 1467--1570, one \texttt{proof} environment of 104 lines). No mathematics is written, repaired or re-derived: statement and proof are the article's own text line for line. Required context retained uncounted: smwuu (lines 200--203); dcfghm (lines 347--353); qsswd (lines 385--408); prop:ultrametric-cauchy (lines 471--478); renjh (lines 666--674); thm:3.1 (lines 815--821). Every definition, display and label of those lines is retained unchanged. Omitted from this package and not redistributed: 1--199, 204--346, 354--384, 409--470, 479--665, 675--814, 822--1453, 1465--1466, 1571--2062. Nothing inside a retained excerpt is omitted or altered: every retained body is delivered exactly as the article's lines are written, so no omission receipt of any kind accompanies this package. Citations: the retained excerpts carry no citation command at all, so no bibliography entry of the article is transcribed and no reference is redistributed. Reference adaptation: none was needed and none was made; every reference inside the delivered unit resolves inside it, so no unresolved reference or citation is produced. Printed slips kept literal: no printed word, symbol, hypothesis or number of the counted statement or of its proof was altered. Layout and class: the article's own class is \texttt{article} with packages including enumitem, lineno, amsfonts, amsmath, amsthm, amscd, amssymb, latexsym, amsbsy, pb-diagram, cite, caption, url, tikz; the wrapper is the lane's pinned \texttt{amsart} editorial wrapper at 10pt on A4, which restates the theorem-like environments the retained text needs and adds the article's own macro definitions that the retained text needs (none), with extra wrapper packages (bm, bbm, dsfont, xspace, cancel, mathtools, enumitem). The delivered numbering is the unit's own. Assets: no retained span contains a figure environment, a graphics-inclusion command, a PDF-page inclusion, a table or a TikZ picture; no image file is copied into this package, the package declares no asset dependency and no image file is redistributed with it. Licence: version arXiv:2510.01667v1 of this article is distributed under the Creative Commons Attribution 4.0 International licence, as recorded in the arXiv licence metadata for this exact version and shown on the abstract page https://arxiv.org/abs/2510.01667v1. A full rights notice is delivered at licenses/arxiv-2510.01667-CC-BY-4.0.txt. Characters: every retained excerpt is pure ASCII, so the delivered engine, XeLaTeX with its default Latin Modern text font, reports no missing character.
\begin{equation}
    \label{smwuu}
D_0(X):=D(X)\setminus \{0\}.
\end{equation}

\begin{lemma}\label{dcfghm}
Let $(X,d)$ be a finite ultrametric space and let $B_1, B_2$ be two different open balls in $(X,d)$. If $B_1 \cap B_2 \neq 0$ holds, then we have either $B_1 \subset B_2$ or $B_2 \subset B_1$. If $B_1$ and $B_2$ are disjoint, then the equality
\[
d(x_1, x_2) = \operatorname{diam}(B_1 \cup B_2)
\]
holds for all $x_1 \in B_1$ and $x_2 \in B_2$.
\end{lemma}

\begin{proposition}\label{qsswd}
Let $(X,d)$ be a discrete metric space.  
Then the following statements are equivalent:
\begin{itemize}[left=10pt]
    \item[(i)] $(X,d)$ is not metrically discrete.
    \item[(ii)] The equality
    \begin{equation}
        \label{llkk}
        \inf D_0(X) =0
        \end{equation}
    holds, where $D_0(X)$ is defined by~\eqref{smwuu}.
    \item[(iii)] There are a sequence $(x_n)_{n\in\mathbb{N}}$ of distinct points of $X$ 
    and a strictly decreasing sequence $(r_n)_{n\in\mathbb{N}}$ of  positive real numbers 
    such that 
    $$
    \lim\limits_{n\to\infty}r_n=0
    $$
    and
$$
\left|B_{r_n}(x_n)\right|=1,\quad \left|B_{kr_n}(x_n)\right|\geq 1
$$
    for each $n\in {\mathbb N}$ and every $k\in (1,\infty).$
\end{itemize}
\end{proposition}

\begin{proposition}\label{prop:ultrametric-cauchy}
Let $(X,d)$ be an ultrametric space. A sequence $(x_n)_{n\in\mathbb{N}}$ of points of $X$ 
is a Cauchy sequence if and only if the limit relation
\begin{equation*}
\lim_{n\to\infty} d(x_n, x_{n+1}) = 0
\end{equation*}
holds.
\end{proposition}

\begin{corollary}\label{renjh}
Let $K=K_{n_1,\ldots,n_k}$ be a complete multipartite finite graph. 
Then the following statements are equivalent:
\begin{itemize}[left=10pt]
    \item[(i)] The equality $k=2$ holds and, in addition, we have  $n_1=n_2=2$.

    \item[(ii)] $K_{n_1,\ldots,n_k}$ is isomorphic to the cycle $C_4$.
\end{itemize} 
\end{corollary}

\begin{theorem}\label{thm:3.1}
Let $(X,\rho)$ be an ultrametric space with $2\leq |X |<\infty$
and let $G_{X,\rho}$ be the diametrical graph of $(X,\rho)$.
Then $G_{X,\rho}$ is a complete $k$-partite finite graph with $k\geq 2$. Conversely, if $G$ is a complete 
$k$-partite finite graph with $k\geq 2$, then there is an ultrametric 
$d \colon V(G) \times V(G) \to \mathbb{R}^+$ such that $G=G_{V,d}$.
\end{theorem}

\begin{lemma}\label{476}
Let an infinite discrete ultrametric space $(X,d)$ satisfy the following conditions:
\begin{itemize}[left=10pt]
    \item[(i)] $(X,d)$ is not metrically discrete.

        \item[(ii)]  $(X,d)$ contains no four-point subspace whose diametrical graph is isomorphic to the cycle $C_4$.

Then $(X,d)$  contains a Cauchy sequence $(x_n)_{n\in \mathbb N}$ of distinct points of $X$.
\end{itemize}

\end{lemma}

\begin{proof}
Since $(X,d)$ is not metrically discrete, Proposition~\ref{qsswd} implies that there are a sequence $(x_n)_{n\in\mathbb{N}}$ of distinct points of $X$ 
    and a strictly decreasing sequence $(r_n)_{n\in\mathbb{N}}$ of  positive real numbers 
    such that 
    \begin{equation}\label{st1}
    \lim\limits_{n\to\infty}r_n=0
        \end{equation}
    and
\begin{equation}
    \label{st2}
\left|B_{r_n}(x_n)\right|=1,\quad \left|B_{kr_n}(x_n)\right|\geq 1
    \end{equation}
    for each $n\in {\mathbb N}$ and every $k\in (1,\infty).$



We claim that $(x_n)_{n\in\mathbb{N}}$ is a Cauchy sequence in $(X,d)$.  
By Proposition~\ref{prop:ultrametric-cauchy} the sequence $(x_n)_{n\in\mathbb{N}}$ is a Cauchy sequence if and only if the limit relation
\begin{equation*}
\lim_{n\to\infty} d(x_n,x_{n+1})=0
\end{equation*}
holds.
The last equality is satisfied if and only if
\begin{equation}\label{xx2}
    \limsup_{n\to\infty} d(x_n, x_{n+1}) = 0.
\end{equation}
Let us assume that \eqref{xx2} is false.  
Then there exist $c_1>0$ and a subsequence $(x_{n_m})_{m\in\mathbb{N}}$ of the sequence $(x_n)_{n\in\mathbb{N}}$ such that
\begin{equation}
    \label{tvhukp}
      d(x_{n_m}, x_{n_m+1}) > c_1
\end{equation}
for every $m\in \mathbb N$.
Now using \eqref{st1} and \eqref{tvhukp} we can find $m_0\in \mathbb N$ such that 
\begin{equation}\label{cc1}
    r_{n_{m_0}+1}< r_{n_{m_0}}<c_1.
\end{equation}
Therefore, there exist $k\in (1,\infty)$ for which the double inequality
\begin{equation}
    \label{cc2}
kr_{n_{m_0}+1}<kr_{n_{m_0}}<c_1
\end{equation}
is satisfied.

For convenience, we introduce the following notation:
\begin{equation}
    \label{cc3}
    x^0 := x_{n_{m_0}}, \quad x^1 := x_{n_{m_0}+1}, \quad
    r^0 := kr_{n_{m_0}}, \quad r^1 := kr_{n_{m_0}+1}.
\end{equation}

Let us consider the open balls $B_{r^0}(x^0)$ and $B_{r^1}(x^1)$.  

Inequality \eqref{tvhukp} with $m=m_0$, double inequality \eqref{cc2}  with $x_m=x_{n_m} = x^0$, and Lemma~\ref{dcfghm} with 
    $B_1 = B_{r^0}(x^0),$ $ B_2 = B_{r^1}(x^1),$
 give us the equality
\begin{equation}
    \label{kkjuuh1}
    B_{r^0}(x^0) \cap B_{r^1}(x^1) = \varnothing.
\end{equation}
The second inequality in \eqref{st2} implies the existence of some points
$
y^0 \in B_{r^0}(x^0),$ and $ y^1 \in B_{r^1}(x^1)
$
which satisfy conditions
\begin{equation}
\label{kkhh4}
y_0 \neq x^0, \qquad y^1 \neq x^1. 
\end{equation}
Now using \eqref{kkjuuh1} and \eqref{kkhh4} we see that the point $x^0,$ $y^0,$ $x^1$ and $y^1$ are pairwise distinct.
Let us denote by $A$ the set $\{x^0, y^0, x^1, y^1\}$ and consider the ultrametric space $(A, d_A)$ with
\[
d_A := d\big|_{A \times A}.
\]

By Theorem \ref{thm:3.1} the diametrical graph $G_{A,d_A}$ is a complete multipartite graph.  
We claim that $G_{A,d_A}$ and $K_{2,2}$ are isomorphic.

To prove it we note that the equalities
\begin{equation}
    \label{865gny}
\operatorname{diam} A=d(x^0, x^1) = d(x^0, y^1) = d(y^0, x^1) = d(y^0, y^1)
\end{equation}
hold by Lemma \ref{dcfghm}.
Moreover, inequality \eqref{tvhukp} with $m=m_0$, double inequality~\eqref{cc1}, 
and equalities~\eqref{cc3} give us the inequalities
\begin{equation}
    \label{vnhgefk}
d(y^0,x^0) <\operatorname{diam} A
\end{equation}
and
\begin{equation}
\label{eftyi88}
d(y^1,x^1) <\operatorname{diam} A. 
\end{equation}
Now using \eqref{865gny}--\eqref{eftyi88} , we see that $G_{A,d_A}$ and $K_{2,2}$ are isomorphic.  

Hence $G_{A,d_A}$ and $C_4$ also are isomorphic by Corollary~\ref{renjh}.  
The last statement contradicts condition $(ii)$.  

Thus $(x_n)_{n \in \mathbb{N}}$ is a Cauchy sequence in $(X,d)$. 


\end{proof}

\end{document}
